\begin{lemma}[Finite-prefix agreement survives finitely many shifts]
If two increasing enumerations agree through their first $n+j$ entries, then
their $j$-fold shifts agree through their first $n$ entries:
\[
\mathsf{PrefixAgree}(n+j,x,y)
\Longrightarrow
\mathsf{PrefixAgree}\bigl(n,\mathsf{PowerShiftIter}(j,x),
  \mathsf{PowerShiftIter}(j,y)\bigr).
\]
\end{lemma}
\begin{proof}
Induct on $j$.  For $j=0$ the assertion is reflexive.  At a successor,
By the definition of \noderef{PrefixAgree}, the first $n$ entries of the shifted sequences are the entries numbered
$1$ through $n$ of the preceding shifts, by the definitions of
\noderef{PowerShiftIter} and \noderef{ShiftMap}.  The hypothesis
\mathsf{PrefixAgree}(n+(j+1),x,y) supplies agreement through the longer
prefix needed for the induction hypothesis at the preceding shift; applying
that hypothesis and then unfolding one more use of \noderef{ShiftMap} gives
agreement through $n$ entries.  This proves the displayed implication.
\end{proof}
